\section{Programação para Inteligência Artificial}

\textbf{OBJETIVO GERAL:}

Desenvolver competências para projetar, implementar e testar algoritmos utilizando a linguagem Python, aplicando fundamentos de programação estruturada, programação orientada a objetos, manipulação de dados e desenvolvimento de aplicações voltadas à Inteligência Artificial.

\textbf{CARGA HORÁRIA:} 120 horas

\textbf{FUNÇÃO:}

Desenvolver sistemas computacionais baseados em Inteligência Artificial, Análise Preditiva e Ecossistemas IoT, atendendo às normas de qualidade corporativas, robustez metodológica, governança de dados e arquitetura segura.

\textbf{SUBFUNÇÕES:}

\begin{itemize}
\item Desenvolver algoritmos utilizando linguagem Python.
\item Estruturar programas utilizando boas práticas de programação.
\item Manipular estruturas de dados para processamento computacional.
\item Desenvolver aplicações orientadas a objetos.
\item Integrar bibliotecas para aplicações em Inteligência Artificial.
\end{itemize}

\textbf{PADRÕES DE DESEMPENHO:}

\begin{itemize}
\item Desenvolver algoritmos corretos e eficientes.
\item Aplicar estruturas de controle de fluxo.
\item Manipular coleções de dados.
\item Modularizar aplicações.
\item Implementar tratamento de exceções.
\item Desenvolver aplicações orientadas a objetos.
\item Utilizar bibliotecas externas.
\item Documentar código seguindo boas práticas.
\end{itemize}

\textbf{CAPACIDADES TÉCNICAS:}

\textbf{Lógica de Programação}

\begin{itemize}
\item Desenvolver algoritmos utilizando raciocínio lógico.
\item Interpretar fluxogramas e pseudocódigo.
\item Resolver problemas computacionais.
\item Estruturar soluções de forma modular.
\end{itemize}

\textbf{Programação em Python}

\begin{itemize}
\item Utilizar sintaxe da linguagem Python.
\item Manipular variáveis, operadores e expressões.
\item Desenvolver estruturas condicionais e de repetição.
\item Implementar funções reutilizáveis.
\end{itemize}

\textbf{Estruturas de Dados}

\begin{itemize}
\item Manipular listas, tuplas, conjuntos e dicionários.
\item Desenvolver algoritmos utilizando coleções.
\item Processar dados estruturados.
\item Organizar informações para aplicações de IA.
\end{itemize}

\textbf{Programação Orientada a Objetos}

\begin{itemize}
\item Desenvolver classes e objetos.
\item Aplicar encapsulamento.
\item Aplicar herança e polimorfismo.
\item Organizar aplicações utilizando orientação a objetos.
\end{itemize}

\textbf{CONHECIMENTOS:}

\textbf{Fundamentos de Programação}

\begin{itemize}
\item Algoritmos
\item Fluxogramas
\item Pseudocódigo
\item Variáveis
\item Tipos de dados
\item Operadores
\end{itemize}

\textbf{Estruturas de Controle}

\begin{itemize}
\item Estruturas condicionais
\item Estruturas de repetição
\item Operadores lógicos
\item Controle de fluxo
\end{itemize}

\textbf{Funções}

\begin{itemize}
\item Declaração de funções
\item Parâmetros
\item Retorno
\item Escopo de variáveis
\item Recursividade
\end{itemize}

\textbf{Estruturas de Dados}

\begin{itemize}
\item Listas
\item Tuplas
\item Dicionários
\item Conjuntos
\item Strings
\end{itemize}

\textbf{Programação Orientada a Objetos}

\begin{itemize}
\item Classes
\item Objetos
\item Métodos
\item Atributos
\item Encapsulamento
\item Herança
\item Polimorfismo
\end{itemize}

\textbf{Tratamento de Exceções}

\begin{itemize}
\item Try
\item Except
\item Finally
\item Raise
\item Depuração
\end{itemize}

\textbf{Manipulação de Arquivos}

\begin{itemize}
\item Leitura de arquivos
\item Escrita de arquivos
\item CSV
\item JSON
\item Organização de dados
\end{itemize}

\textbf{Bibliotecas para IA}

\begin{itemize}
\item NumPy
\item Pandas
\item Matplotlib
\item Requests
\item Ambientes virtuais
\end{itemize}

\textbf{Boas Práticas}

\begin{itemize}
\item PEP 8
\item Documentação
\item Modularização
\item Reutilização de código
\item Versionamento
\end{itemize}

\textbf{CAPACIDADES SOCIOEMOCIONAIS:}

\begin{itemize}
\item Demonstrar organização na construção de algoritmos.
\item Desenvolver raciocínio lógico para resolução de problemas.
\item Avaliar criticamente soluções computacionais.
\item Trabalhar colaborativamente utilizando controle de versões.
\item Produzir documentação técnica de qualidade.
\item Compartilhar conhecimento com equipes de desenvolvimento.
\item Demonstrar autonomia na busca por soluções.
\item Agir com responsabilidade na construção de aplicações computacionais.
\end{itemize}

\textbf{AMBIENTES PEDAGÓGICOS:}

\begin{itemize}
\item Laboratório de Informática.
\item Laboratório de Desenvolvimento de Software.
\item Ambiente Virtual de Aprendizagem.
\item Ambiente Cloud para execução de aplicações.
\end{itemize}

\textbf{EQUIPAMENTOS E RECURSOS:}

\textbf{Hardware}

\begin{itemize}
\item Computadores com acesso à Internet.
\item Projetor multimídia.
\item Ambiente para desenvolvimento colaborativo.
\end{itemize}

\textbf{Software}

\begin{itemize}
\item Python
\item VS Code
\item Git
\item GitHub Desktop
\item Jupyter Notebook
\item Google Colab
\item Docker (introdução)
\end{itemize}
